Sampling-based semantic consistency --- computing NLI agreement across multiple stochastic samples --- is a widely adopted training-free hallucination detector for black-box LLMs, validated on open-ended generation tasks with base language models. We evaluate this paradigm on Llama-3-8B-Instruct, a representative RLHF-fine-tuned instruction model, using NLI-based (SMC-NLI) and embedding-based (SMC-Embed) consistency scores on 1,000 balanced HaluEval QA questions. Both metrics achieve chance-level AUROC (SMC-NLI: 0.4933; SMC-Embed: 0.4859), with implementation correctness independently verified by 14 unit tests and a mechanism verification step. Analysis reveals the failure is distributional: mean SMC-NLI scores for correctly-labeled (0.6236) and hallucinated-labeled (0.6299) questions differ by only 0.006, indicating that Llama-3-8B-Instruct produces consistent outputs for both correct and wrong beliefs --- a ``systematic confabulation'' regime induced by RLHF fine-tuning that renders consistency-based signals non-discriminative. We introduce the stochastic hallucination vs.\ systematic confabulation distinction as a necessary regime prerequisite for sampling-based detectors, provide empirical evidence that instruction-tuned models on structured factual QA violate this prerequisite, and release validated SMC infrastructure for future evaluation on tasks where the stochastic regime may hold.
